\section{Discussion and limitations}
\label{sec:discussion}

\paragraph{Measured versus literal.} The compute and memory wins of
\S\ref{sec:sharedkv} are \emph{measured}: the candidate-pool probe reports how many
candidates a real index would score but still computes $q\!\cdot\!K$ densely to
obtain the recall number, and the quantization probe fake-quantizes in fp32. The
results establish that recall survives the pool and 4 bits; realising the actual
speedup and byte savings needs a genuine token-hash / LSH kernel and int4 storage.
That is engineering, not a further scientific claim, but it is not done here.

\paragraph{Eval-time read path.} The pool and quantization are applied at
evaluation on models trained with the plain dense, full-precision read. Evaluation
showed no gap, so we did not need to train with them; whether training \emph{with}
the pooled/quantized read helps at harder settings is open.

\paragraph{Curriculum is an optimisation, not a requirement.} We initially reported
that the warm-start curriculum was needed for full recall on real language. Running
the cold multi-layer model to a larger budget overturns this: it reaches
\recall{99.7} at 16 facts and \recall{98.0} at 32 on its own, through a delayed
phase transition, and the curriculum reaches the same endpoint in roughly half the
steps (\S\ref{sec:placement}). So the curriculum buys training efficiency, not
capability. This is a caution as much as a result: a plateau before a delayed
transition is easy to mistake for a hard bottleneck, and the ``embedding
competition'' wall we posited was partly an under-training artifact. How the
transition step scales with fact count, model size, and seed---and whether a
principled schedule beats both cold-long and the hand-built ladder---remains open.

\paragraph{The exact-key constraint.} The layer-0 exact-token key is both the
enabling property and a limitation: it makes retrieval exact and cheap but forbids
contextual keys and pattern separation, and top-$1$ sparse reading works precisely
because the answer slot dominates an exact-match ranking. On tasks with genuinely
fuzzy keys a larger pool (the LSH index, $k>1$) is needed, and the clean $O(1)$
token hash no longer applies.

\paragraph{Scope.} Some placement and breaking-point results are on a synthetic
associative-recall task; the flagship recall, long-context, sparse-read, and
shared-KV results are on real PG-19 language. The model is small (6 layers,
$256$-d, single head); we have not tested scale, and the consolidation/sleep
machinery is validated only on the synthetic store.
